\documentclass[10pt]{article}

\usepackage[utf8]{inputenc}

\usepackage[T1]{fontenc}

\usepackage{amsmath}

\usepackage{amsfonts}

\usepackage{amssymb}

\usepackage[version=4]{mhchem}

\usepackage{stmaryrd}



\begin{document}

С л е д с т в и е. Абсолютными инвариантами фазового потока являются и формы $\omega^{1}, \omega^{2}, \ldots, \omega^{n}$.



Последняя из этих форм является фазо в ы м объе мом. Сле д с т в е (Лиувилля теорема о фазовом объеме). Фазовый поток сохраняет объем в фазовом пространстве.



Если гамильтонова система имеет первые интегралы $F_{1}$, $\ldots, F_{m}$, то сужение фазового потока $g^{t}$ на неособое инвариантное многообразие $M_{f}=\left\{F_{i}=f_{i}, i=1, \ldots, m\right\}$ сохраняет нек-рую меру с гладкой положительной плотностью: в качестве такой меры можно взять форму $d \sigma / V_{m}$, где $d \sigma-$ элемент объема $M_{f}$ как вложенного подмногообразия, а $V_{m}-m$-мерный объем параллелепипеда со сторонами grad $F_{1}, \ldots, \operatorname{grad} F_{m}$. Эта форма определяется лишь симплектич. структурой и не зависит от выбора метрики на $M$.



М. Э. Казарян.



ИНТЕГРАЛЬНЫЙ КОСИНУС - см. Интегральные функции. ИНТЕГРАЛЬНЫЙ ЛОГАРИФМ - сМ. ИНтегральНЫе функции.



ИНТЕГPАЛЬНЫЙ OПEPATOP - обобщение понятия матрицы на бесконечномерный случай.



Матрица $K_{i j}$ отображает векторы $x_{j}$ из векторного пространства $X$ в векторы $y_{i}=K_{i j} x_{j}$ пространства $Y$. Простейший линейный И.о., определяемый равенством



\[

y(t)=\int_{S} K(t, s) x(s) d s,

\]



отображает функции $x(s)$ из функционального пространства $X$ (области определения) в функции $y(t)$ из функционального пространства $Y$ (область значений); функция $K(t, s)$ называется я д р о м И.о. Чаще всего рассматривают И. о. на функциональных пространствах $C(S)$ (непрерывных на замкнутом множестве $S$ функций) и $L^{p}(S)$ (интегрируемых на $S$ со степенью $p$ функций).



Среди И.о. наиболее изучены (вполне непрерывные) фредгольмовы операторы. Ядро $K$ при этом называется фредгольмовым ядром. Напр., для И.о., действующего в $C(S)$, ядро $K$ фредгольмово, если функция $K(t, s)$ непрерывна в квадрате $S \times S$. Для И. о. в $L^{2}(S)$ ядро фредгольмово, если выполнено неравенство:



\[

\int_{S} \int_{S}|K(t, s)|^{2} d t d s<\infty .

\]



Важным частным случаем фредгольмова оператора является оператор Гильберта-Шмидта (см. Интегральное уравнение). Встречаются И.о. с полярным ядром (со слабой особенностью):



\[

K(t, s)=B(t, s)|t-s|^{-m}, 0<m<n,

\]



где $|t-s|-$ расстояние между точками $s$ и $t n$-мерного пространства. Для функший из $C(S)$ И.о. с полярным ядром будет фредгольмовым, если функция $B(t, s)$ непрерывна на $S \times S$; если $B(t, s)$ ограничена всюду в квадрате $S \times S$ и



\[

\int_{S} \int_{S}|B(s, t)|^{2} d s d t<\infty,

\]



то И. о. с полярным ядром фредгольмов на $L^{2}(S)$.



В математич. физике применяют различные типы И.о., возникающих при интегральных преобразованиях.



Juт.: [1] В лад и м и р о в В. С., Уравнения математической физики, 5 изд., М., 1988; [2] Интегральные уравнения, М., 1968; [3] Рихтмай е р Р., Принципы современной математической физики, пер. с англ., М., 1982. С. В. Молодиов.



ИНТЕГРАЛЬНЫЙ СИНУС - см. Интегральные функции.



UHTERPИPOBAHИE по антикоммутирующим пе ре ме н ным - интегрирование по образующим $x_{1}, \ldots, x_{n}$, $x_{1}^{+}, \ldots, x_{n}^{+}$алгебры Грассмана, определяемое следующим образом. Вводятся символы $d x$, подчиненные соотношениям антикоммутативности:



\[

\left.\begin{array}{c}

\left\{d x_{i}, d x_{j}\right\}=\left\{d x_{i}, d x_{j}^{+}\right\}=\left\{d x_{i}^{+}, d x_{j}^{+}\right\}=0, \\

\left\{x_{i}, d x_{j}\right\}=\left\{x_{i}, d x_{j}^{+}\right\}=\left\{x_{i}^{+}, d x_{j}\right\}=\left\{x_{i}^{+}, d x_{j}^{+}\right\}=0,

\end{array}\right\}

\]



где $\{a, b\}=a b+b a$. Определяются однократные интегралы



\[

\left.\begin{array}{rl}

\int d x_{i} & =\int d x_{i}^{+}=0, \\

\int x_{i} d x_{i} & =\int x_{i}^{+} d x_{i}^{+}=1 .

\end{array}\right\}

\]



Кратные интегралы понимают как повторные. Формулы (1) и (2) однозначно определяют интеграл от любого одночлена $x_{n}^{\alpha_{n}} \cdot \ldots \cdot x_{1}^{\alpha_{1}} \cdot x_{n}^{+\beta_{n}} \ldots \cdot x_{1}^{+\beta_{1}}$. На произвольные элементы грассмановой алгебры



\[

\begin{aligned}

f\left(x, x^{+}\right) & =\sum_{\alpha i, \beta_{i}=0,1} C\left(\alpha_{1}, \ldots, \alpha_{n} ; \beta_{1}, \ldots, \beta_{n}\right) \times \\

& \times\left(x_{n}\right)^{\alpha_{n}} \ldots\left(x_{1}\right)^{\alpha_{1}}\left(x_{n}^{+}\right)^{\beta_{n}} \ldots\left(x_{1}^{+}\right)^{\beta_{1}}

\end{aligned}

\]



интеграл распространяется по линейности.



Такие правила впервые последовательно были введены Ф. А. Березиным (см. [1], [2]). И. по антикоммутирующим переменным естественным образом возникает при квантовании ферми-полей в формализме континуального интегрирования (см. [3], [4]).



Для произвольного элемента $f\left(x, x^{+}\right)$



\[

\begin{gathered}

\int f\left(x, x^{+}\right) d x^{+} d x \equiv \\

\equiv \int f\left(x, x^{+}\right) \prod_{i=1}^{n} d x_{i}^{+} d x_{i}=(-1)^{n-1} C(1, \ldots, 1 ; 1, \ldots, 1),

\end{gathered}

\]



где $C$ - коэффициенты из записи элемента $f\left(x, x^{+}\right)$в форме (3).



Для интегралов на грассмановой алгебре справедлива формула замены переменных. Ограничиваясь случаем, когда замена линейна, то есть



\[

\begin{aligned}

x_{i}=\sum_{j} a_{i j} x_{j}^{\prime}, & d x_{i}=\sum_{j} a_{i j}^{-1} d x_{j}^{\prime}, \\

x_{i}^{+}=\sum_{j} b_{i j} x_{j}^{\prime+}, & d x_{i}^{+}=\sum_{j} b_{i j}^{-1} d x_{j}^{+},

\end{aligned}

\]



где $a_{i j}$ - элементы матрицы $A, a_{i j}^{-1}$ - элементы обратной матрицы $A^{-1}$, получим



\[

\int f\left(x, x^{+}\right) d x^{+} d x=\operatorname{det}\left(A^{-1} B^{-1}\right) \int f\left(a x^{\prime}, b x^{+}\right) d x^{\prime+} d x^{\prime} .

\]



Для гауссовых интегралов имеются формулы (см. [2]-[4]):



\[

\begin{gathered}

\int \exp \left(-x^{+} A x+\eta^{+} x+x^{+} \eta\right)=\operatorname{det} A \cdot \exp \left(\eta^{+} A^{-1} \eta\right), \\

\int \exp (-x B x) \prod_{i=1}^{n} d x_{i}=\operatorname{det}^{1 / 2} B,

\end{gathered}

\]



где



\[

\begin{gathered}

x A x=\sum x_{i}^{+} a_{i j} x_{j} ; \quad x B x=\sum x_{i} b_{i j} x_{j} \quad\left(b_{i j}=-b_{j i}\right), \\

\eta^{+} x=\sum \eta_{i}^{+} x_{i}, \quad x^{+} \eta=\sum x_{i}^{+} \eta_{i}, \\

\left\{\eta_{i}, \eta_{j}\right\}=\left\{\eta_{i}^{+}, \eta_{j}\right\}=\left\{\eta_{i}^{+}, \eta_{j}^{+}\right\}=0 .

\end{gathered}

\]



Лum.: [1] Бе рез и н Ф.А., «Докл. АН СССР», 1961, т. 137, № 2, c. 311-14; [2] е го же, Метод вторичного квантования, 2 изд., М., 1986; [3] В а с и л ь е в А. Н., Функциональные методы в квантовой теории поля и статистике, Л., 1976; [4] По по в В. Н., Континуальные интегралы в квантовой теории поля и статистической физике, $\begin{array}{ll}\text { М., 1976. } & \text { C. A. Федотов. }\end{array}$ ИНТЕГPИРОВАНИЯ МЕТОД По быстрым и медле нны м пе ре менным - метод вычисления длинноволновых асимптотик функций Грина в квантовой теории поля и статистической физике в формализме функщионального интеграла, аналог метода усреднения по быстрым движениям в нелинейной механике (см. [1]). Если в системе присутствуют кванты сколь угодно малой энергии, члены стандартной теории возмущений имеют так наз. инфракрасные сингулярности, поэтому желательна перестройка теории возмущений, к-рую и позволяет осуществить И. м. по быстрым и медленным переменным (см. [2], [3]).





\end{document}